% ************************** Thesis Abstract *****************************
% Use `abstract' as an option in the document class to print only the titlepage and the abstract.
% \renewcommand{\abstract}{myuseristhis}
\begin{spanishabstract}

Los sistemas de bicicleta compartida (SBC) son una opción de movilidad accesible y de bajas emisiones que puede sustituir trayectos urbanos cortos realizados en coche y mejorar la conexión con el transporte público. A medida que estos sistemas han evolucionado hacia servicios cada vez más grandes, su rendimiento depende cada vez más de decisiones tomadas a lo largo de todo el ciclo de vida, desde la ubicación de las estaciones hasta la operativa diaria. Esta tesis desarrolla métodos basados en datos para apoyar estas tareas, utilizando Barcelona y su sistema con estaciones fijas, \textit{Bicing}, como caso de estudio.

En primer lugar, la tesis aborda el diseño de la red de estaciones. Los enfoques de optimización existentes a menudo no permiten cuantificar de forma realista la accesibilidad del sistema, ni explorar de manera transparente los equilibrios entre distintos objetivos de planificación (p. ej., maximizar la cobertura de la demanda o la equidad social), ni controlar explícitamente la disposición de las estaciones, que suele emerger como un resultado indirecto de los procesos de optimización. Esta tesis propone un método de planificación flexible que combina una puntuación multicriterio modular con métricas de accesibilidad basadas en la red viaria dirigida y ajustadas a la topografía, y utiliza un algoritmo genético para optimizar la disposición de las estaciones, controlando explícitamente la proximidad entre ellas y la accesibilidad global del sistema. Aplicado a Barcelona, el método identifica configuraciones cercanos al óptimo para distintos escenarios y muestra que la accesibilidad ajustada por la topografía contribuye a corregir déficits de servicio en zonas con pendiente. Asimismo, el método permite identificar nuevas estaciones que complementan una red existente.

En cuanto a la operativa diaria, la tesis comienza caracterizando los patrones de movilidad de \textit{Bicing}, que sirven de base para dos casos de uso predictivos. El primero tiene como objetivo maximizar la disponibilidad de bicicletas eléctricas manteniendo las baterías cargadas, sin disponer de información exacta sobre las rutas ni de mediciones de batería en tránsito. La tesis propone un método de predicción conjunto que combina un modelo de movilidad basado en cadenas de Markov con un modelo de batería fundamentado en principios físicos para simular trayectorias y niveles de carga. Este enfoque permite obtener estimaciones precisas de la batería y predicciones fiables a corto plazo a nivel de estación, facilitando la aplicación de estrategias de carga proactivas.

El segundo caso de uso también tiene como objetivo aumentar la disponibilidad, pero reduciendo el tiempo de inactividad de las bicicletas mediante la anticipación de fallos en sus componentes. La tesis desarrolla modelos de mantenimiento predictivo para tres componentes clave (pastillas de freno, radios de la rueda y cadenas) basados en análisis de supervivencia y aprendizaje automático, utilizando registros de mantenimiento, indicadores históricos de movilidad y variables meteorológicas para estimar el riesgo de fallo. Los modelos muestran un rendimiento predictivo preciso y consistente, y el análisis de interpretabilidad indica que un uso más intensivo acorta la vida útil de los componentes, en particular una mayor velocidad media y la propulsión eléctrica.

En conjunto, la tesis demuestra que la planificación y la operativa de los SBC con estaciones fijas pueden mejorarse mediante herramientas computacionales basadas en datos, incluso en contextos con limitaciones de datos. Los métodos propuestos permiten pasar de una planificación heurística y una gestión reactiva a estrategias anticipativas basadas en datos. Además, son transferibles a otras ciudades con datos espaciales y operativos comparables.

\textbf{Palabras clave:} bicicleta compartida, movilidad eléctrica, optimización espacial, métodos basados en datos, ubicación de estaciones, predicción de batería, mantenimiento predictivo, Barcelona, Bicing.

\end{spanishabstract}
